\chapter{Background}\label{chap:background-foundations}

\section{Emotion Recognition in Learning}

Emotion recognition means identifying what a learner is feeling from observable signals, such as facial expressions in video frames \cite{lek_academic_2023,florestiyanto_emotion_2024}. In this thesis, the task is Facial Expression Recognition (FER), which is the use of computer vision to detect emotion-related facial patterns automatically \cite{lek_academic_2023}. FER is useful in learning systems because it can capture signals continuously without asking the learner to stop and fill in surveys \cite{vistorte_integrating_2024}.

Educational research usually focuses on academic emotions, not only basic emotions like happiness or fear \cite{pekrun_control-value_2006,dmello_language_2012}. Academic emotions are feelings linked to studying and solving tasks, such as engagement, confusion, frustration, and boredom \cite{dmello_language_2012}. These states matter because they affect attention, effort, and persistence during learning \cite{pekrun_control-value_2006,tyng_influences_2017}.

The DAiSEE dataset is an educational video dataset where clips are labeled for engagement, boredom, confusion, and frustration \cite{gupta_daisee_2016}. A dataset is a structured collection of examples used to train and test machine learning models \cite{lek_academic_2023}. In this thesis, DAiSEE-style labels are used as the emotion targets for the FER module \cite{gupta_daisee_2016}.

\section{Deep Learning Components Used in This Thesis}

Deep learning is a machine learning approach where a model learns many layers of representations directly from data \cite{aly_revolutionizing_2025}. A neural network is the model structure used in deep learning; it maps input values to output predictions through connected layers \cite{gupta_edfa_2023}. In FER, deep learning helps the system learn complex facial patterns that are hard to write as manual rules \cite{lek_academic_2023}.

A Convolutional Neural Network (CNN) is a neural network designed for images \cite{aly_revolutionizing_2025}. It uses convolution filters to detect local patterns, such as edges, textures, and face parts, then combines them into higher-level features \cite{gupta_edfa_2023}. In this thesis, the CNN backbone is ResNet50, which is a 50-layer residual network widely used as an image feature extractor \cite{he_resnet_2016,asaju_adaptive_2022,aly_revolutionizing_2025}.

A feature extractor is the part of a model that converts raw images into compact numeric representations called features \cite{komaravalli_detecting_2023}. A backbone is the main feature extractor network used before the final prediction layers \cite{aly_revolutionizing_2025}. Here, ResNet50 is used as the backbone to extract frame-level facial features before temporal modeling \cite{asaju_affects_2021,asaju_adaptive_2022}.

Transfer learning means starting from a model that was pre-trained on a large dataset, then fine-tuning it for the target task \cite{komaravalli_detecting_2023}. Pre-trained means the model has already learned generic visual patterns before this project-specific training step \cite{aly_revolutionizing_2025}. This approach is useful when educational FER datasets are smaller than general image datasets \cite{lek_academic_2023}.

Some reviewed studies also use other backbone families. Xception and EfficientNet are CNN architectures designed to improve feature quality and efficiency \cite{komaravalli_detecting_2023,shiri_detection_2024}. A Vision Transformer (ViT) is a model that applies transformer attention to image patches instead of convolution filters \cite{tulsani_realtime_2025,su_leveraging_2024}. A transformer is a sequence model that uses self-attention to model dependencies between elements in parallel \cite{su_leveraging_2024,liu_dynamic_nodate}.

InceptionNet is another CNN family that uses multi-branch convolution blocks with different kernel sizes in parallel \cite{gupta_daisee_2016}. It is often used as a benchmark backbone in FER studies \cite{gupta_daisee_2016}.

Classical components also appear in hybrid FER pipelines. A Support Vector Machine (SVM) is a classifier that separates classes by finding a maximum-margin boundary in feature space \cite{rao_recognition_2021,turkes_enhanced_2025}. PCA and SVD are dimensionality reduction methods that compress features while keeping major information patterns \cite{santoni_convolutional_2023}. SMOTE is an oversampling method that creates synthetic minority examples to reduce class imbalance \cite{santoni_convolutional_2023}.

Some models use specialized feature terms. A GCN, or Graph Convolutional Network, learns from graph-structured inputs such as landmark graphs \cite{rianto_engagement_2025}. Action Units are localized facial muscle movement codes used in facial analysis \cite{santoni_convolutional_2023,rianto_engagement_2025}. LDAM is an imbalance-aware loss that increases class margins for better minority-class separation \cite{rianto_engagement_2025}.

\section{Temporal Modeling in Video}

Temporal modeling means learning how information changes over time, not only what appears in one single frame \cite{asaju_adaptive_2022,mehta_three-dimensional_2022}. In educational FER, many states build gradually, so a short frame sequence often gives better evidence than one image \cite{abedi_improving_2021}. For this reason, this thesis models emotion from sampled video frames, not isolated still images.

Frame sampling means selecting a fixed number of frames from each video clip for processing \cite{abedi_improving_2021}. This makes training stable because each sample has the same temporal length \cite{mehta_three-dimensional_2022}. It also reduces computation compared with using all frames in long clips \cite{asaju_adaptive_2022}.

An LSTM (Long Short-Term Memory) is a recurrent neural network that stores useful past information while processing sequences step by step \cite{hochreiter_lstm_1997,leong_deep_2020}. A BiLSTM (Bidirectional LSTM) reads the sequence in both forward and backward directions, so each time step can use past and future context \cite{asaju_adaptive_2022,tulsani_realtime_2025}. In this thesis, BiLSTM is used after the ResNet50 features to capture temporal emotion patterns in the video clip.

A 3D CNN uses 3D convolutions over height, width, and time together, so it learns spatial and temporal patterns in one operation \cite{mehta_three-dimensional_2022}. C3D is an early 3D CNN design used for short video clips \cite{gupta_daisee_2016}. LRCN means Long-term Recurrent Convolutional Network, which combines CNN frame features with a recurrent temporal module \cite{gupta_daisee_2016}.

\section{Attention Mechanisms}

An attention mechanism is a learnable weighting method that tells the model which parts of its input are more important for the current prediction \cite{vaswani_attention_2017,mehta_three-dimensional_2022,su_leveraging_2024}. Instead of treating all features equally, attention gives larger weights to more informative signals \cite{liu_dynamic_nodate}. This usually improves focus and reduces noise in difficult real-world data \cite{aly_revolutionizing_2025}.

SE attention in this thesis refers to a squeeze-and-excitation style block that recalibrates feature channels by learning channel-wise importance weights \cite{hu_squeeze_2018,aly_revolutionizing_2025}. Channel means one dimension of the feature representation produced by the backbone \cite{liu_dynamic_nodate}. This helps the model strengthen useful facial cues and weaken less relevant channels before temporal encoding.

Temporal attention means applying attention over time steps in a sequence \cite{mehta_three-dimensional_2022}. The model learns which frames in the clip carry the strongest emotional evidence and gives them more influence in the final video representation \cite{su_leveraging_2024,liu_dynamic_nodate}. In this thesis, temporal attention is used after BiLSTM to produce the final sequence feature for classification.

CBAM, or Convolutional Block Attention Module, is an attention block that combines channel and spatial attention in CNN features \cite{aly_revolutionizing_2025}. A Temporal Convolutional Network (TCN) is a sequence model that uses 1D temporal convolutions instead of recurrent units to capture time dependencies \cite{abedi_improving_2021,aly_revolutionizing_2025}. MTCNN is a cascaded face detector that finds and aligns face regions before expression classification \cite{tulsani_realtime_2025}.

DTFormer means Differential Temporal Transformer, which is a transformer-style temporal module designed to model fine-grained frame-to-frame emotional change \cite{liu_dynamic_nodate}. In reviewed FER work, it is combined with adaptive global attention for dynamic learner expression recognition \cite{liu_dynamic_nodate}.

\section{Multilabel Classification and Class Imbalance}

Multilabel classification means one input sample can have more than one active label at the same time \cite{lek_academic_2023}. This is different from single-label classification, where each sample belongs to exactly one class \cite{gupta_daisee_2016}. In practice, a learner clip can show mixed academic affect, so multilabel output is used in this thesis.

A sigmoid activation is a function that converts each output logit into a value between 0 and 1, which can be interpreted as per-label probability \cite{mehta_three-dimensional_2022}. Because sigmoid is applied per label, each label decision is independent from the others \cite{rianto_engagement_2025}. This is why sigmoid is commonly used in multilabel settings.

A loss function is the training objective that measures the difference between model predictions and true labels \cite{abedi_improving_2021}. Weighted loss means giving more weight to under-represented classes so the model does not ignore them during optimization \cite{mehta_three-dimensional_2022,rianto_engagement_2025}. In this thesis, weighted binary cross-entropy is used to handle label imbalance.

Controlled oversampling means increasing the frequency of minority patterns in training data in a planned way, so learning is less biased toward majority classes \cite{santoni_convolutional_2023}. Oversampling does not change the model structure; it changes what the model sees during training \cite{abedi_improving_2021}. This thesis uses controlled oversampling to improve balance across emotion labels.

Threshold tuning means selecting decision cutoffs for each label instead of using one fixed default threshold for all labels \cite{rianto_engagement_2025}. After training, the threshold for each label is adjusted to improve validation performance, especially F1 \cite{mehta_three-dimensional_2022}. This is important in imbalanced multilabel tasks because one threshold rarely fits all labels equally well.

\section{Evaluation Metrics for Multilabel FER}

F1 score combines precision and recall into one value \cite{rianto_engagement_2025}. Precision measures how many predicted positives are correct, and recall measures how many real positives are found \cite{fu_integrating_2025}. F1 is useful when class frequencies are imbalanced, which is common in educational emotion data \cite{mehta_three-dimensional_2022}.

Macro F1 means F1 is computed for each label separately, then averaged equally across labels \cite{rianto_engagement_2025}. This metric treats minority and majority labels with the same importance \cite{mehta_three-dimensional_2022}. It is useful when we want fair performance across all emotion labels.

Micro F1 pools decisions from all labels before computing one global F1 value \cite{fu_integrating_2025}. This metric gives more influence to frequent labels because they contribute more total decisions \cite{rianto_engagement_2025}. Micro F1 is useful for overall system-level performance.

Weighted F1 averages per-label F1 using label frequency as weights \cite{fu_integrating_2025}. It balances between macro and micro behavior by keeping per-label structure while still considering prevalence \cite{rianto_engagement_2025}. This helps summarize practical performance on imbalanced data.

Label accuracy means accuracy computed separately for each label \cite{mehta_three-dimensional_2022}. It shows how well each emotion label is recognized, which is important because some labels are harder than others \cite{lek_academic_2023}. In this thesis, label-level reporting is used together with F1 metrics for a complete evaluation view.

MSE, or Mean Squared Error, is the average squared difference between predicted and target numeric values \cite{mehta_three-dimensional_2022}. It is used when the output is treated as a regression value, such as continuous engagement level prediction \cite{su_leveraging_2024}. Top-1 accuracy means the prediction is counted correct when the model's highest-score class matches the true class \cite{gupta_daisee_2016}.

\section{Adaptive Learning and Intelligent Tutoring Systems}\label{sec:bg:adaptive}

Adaptive learning systems are educational technologies that personalize instruction along several dimensions---content sequencing, pacing, difficulty, hints, and feedback---using data collected during the learning process \cite{gligorea_adaptive_2023,strielkowski_span_2025,zerkouk_comprehensive_2025}. Their defining property is a closed-loop interaction between learner and machine: the system observes the learner's behaviour and outcomes, updates an internal estimate of the learner's state, and selects the next instructional action accordingly \cite{gligorea_adaptive_2023,zerkouk_comprehensive_2025}. Conceptually, this places adaptive learning within the broader class of \emph{sequential decision-making} problems, where each instructional choice has both immediate and long-term effects on learning outcomes \cite{kabudi_ai-enabled_2021}. The adaptation component of this thesis adopts exactly this view: at every interaction step, the tutor must decide an action that is appropriate for the learner's current cognitive and affective condition.

\subsection{Intelligent Tutoring Systems}\label{sec:bg:its}

An Intelligent Tutoring System (ITS) is a computer-based learning environment that delivers individualized tutoring by reasoning over an internal learner model, a domain model, and a pedagogical model \cite{zerkouk_comprehensive_2025,kabudi_ai-enabled_2021}. Classical ITS architectures rely on hand-crafted pedagogical rules to choose hints, examples, or remediation, which limits their ability to optimize long-term learning trajectories across diverse learner populations \cite{guzman_developing_2025}. Recent ITS research replaces these rigid rules with data-driven pedagogical policies, using machine learning to decide when to give hints, change difficulty, or schedule review \cite{axak_adaptive_nodate,deshmukh_developing_2025,guzman_developing_2025,meng_self-evolving_2025}. This shift connects ITS design directly to decision-making frameworks studied in machine learning, in particular reinforcement learning.

\subsection{Knowledge Tracing and Deep Knowledge Tracing}\label{sec:bg:kt}

Knowledge tracing (KT) is the task of estimating a learner's evolving mastery of skills or concepts from observed interactions, typically sequences of exercises tagged with whether the learner answered correctly \cite{urbaite_adaptive_2026,kuling_ken_2025}. Accurate knowledge tracing is a prerequisite for principled adaptive instruction, because the next pedagogical decision depends on what the learner currently knows and is ready to learn next \cite{kuling_ken_2025,urbaite_adaptive_2026}.

Deep Knowledge Tracing (DKT) refers to the family of neural-network-based knowledge tracing models, in which recurrent or memory-augmented networks learn rich, continuous representations of learner mastery directly from raw interaction sequences \cite{urbaite_adaptive_2026,kuling_ken_2025,guzman_developing_2025}. Compared with classical knowledge-tracing models, deep knowledge tracing captures temporal dependencies between exercises more flexibly and can be combined with downstream decision-making components such as exercise recommenders or reinforcement-learning-based tutors \cite{fu_integrating_2025,guzman_developing_2025,meng_self-evolving_2025}.

\section{Reinforcement Learning}\label{sec:bg:rl}

The closed-loop, sequential nature of adaptive instruction motivates reinforcement learning (RL) as a unifying decision-making framework. RL studies how an agent can learn to act in an environment so as to maximize expected long-term reward \cite{sutton_rl_book_2018}. The standard formalism is the Markov Decision Process (MDP), specified by a state space, an action space, a transition function, a reward function, and a discount factor \cite{sutton_rl_book_2018}. At each step the agent observes a state, selects an action according to its policy, receives a scalar reward, and transitions to a new state; the learning objective is a policy that maximizes the expected discounted sum of future rewards \cite{sutton_rl_book_2018}.

This formulation maps almost directly onto adaptive learning: the state encodes a learner's cognitive and affective condition, the action is a pedagogical decision such as selecting an explanation strategy or the difficulty of the next exercise, and the reward measures progress toward learning objectives \cite{axak_adaptive_nodate,perez_emotions_2024,gong_design_2025}. This correspondence is what makes RL a natural choice for tutoring systems that must balance immediate engagement with long-term mastery \cite{axak_adaptive_nodate,deshmukh_developing_2025,fu_integrating_2025}. In this thesis the same view is adopted: the adaptation module is treated as an RL agent whose state combines emotion-aware signals from the FER module with task-progress information.

\subsection{Value-Based Methods and Q-Learning}\label{sec:bg:qlearning}

A central object in many RL algorithms is the action-value function $Q^{\pi}(s,a)$, defined as the expected discounted return obtained by taking action $a$ in state $s$ and following policy $\pi$ thereafter \cite{sutton_rl_book_2018}. The optimal action-value function $Q^{\star}$ satisfies the Bellman optimality equation, and an optimal policy can be obtained by acting greedily with respect to $Q^{\star}$ \cite{sutton_rl_book_2018}. \emph{Q-learning} is an off-policy temporal-difference algorithm that iteratively updates an estimate of $Q^{\star}$ from observed transitions, without requiring a model of the environment \cite{sutton_rl_book_2018}. In its tabular form, Q-learning is theoretically well understood but is restricted to small, discrete state spaces \cite{sutton_rl_book_2018}.

\subsection{Deep Q-Networks}\label{sec:bg:dqn}

The Deep Q-Network (DQN) algorithm extends Q-learning to high-dimensional state spaces by representing the action-value function with a deep neural network \cite{mnih_dqn_2015}. Two algorithmic innovations make deep Q-learning practical. First, an \emph{experience replay} buffer stores past transitions and samples mini-batches from it, which decorrelates consecutive updates and improves data efficiency \cite{mnih_dqn_2015}. Second, a periodically updated \emph{target network} provides stable bootstrapping targets, mitigating the divergence often observed when value-based RL is combined with non-linear function approximation \cite{mnih_dqn_2015}.

DQN and its variants are now common building blocks in adaptive-learning research, where the same ingredients are used to train dynamic instructional policies \cite{axak_adaptive_nodate,sayed_ai-based_2023,kong_real-time_2025,gong_design_2025,govea_implementation_2024}. In this thesis the adaptation module follows the same recipe: a deep Q-network, an experience replay buffer, and a target network are trained on a state representation that combines emotion-aware features with task-progress signals.

\section{Simulation-Based Evaluation in RL Tutoring Systems}\label{sec:bg:simulation}

Reinforcement learning is fundamentally a trial-and-error procedure: an agent must explore many state--action sequences before a useful policy emerges \cite{sutton_rl_book_2018}. Carrying out this exploration directly with human learners is generally impractical, because early policies are by definition under-trained, the sample requirements of RL are large, and any policy update during deployment may inadvertently degrade the learning experience \cite{azoulay_adaptive_2021,axak_adaptive_nodate}. For these reasons, RL-based tutoring research typically trains and evaluates policies inside a learner simulator before any human study \cite{azoulay_adaptive_2021,axak_adaptive_nodate,deshmukh_developing_2025,weitekamp_tutorgym_2025}.

A learner simulator is an executable model of student behaviour that, given the current state and a tutor action, returns a next state and a reward signal \cite{azoulay_adaptive_2021,weitekamp_tutorgym_2025,axak_adaptive_nodate}. Simulators may incorporate parameters of cognition, such as mastery and response accuracy; behaviour, such as help-seeking and response time; or affect, such as engagement and frustration \cite{axak_adaptive_nodate,perez_emotions_2024}. In the methodology of this thesis the simulator plays three complementary roles. First, it makes policy training feasible by providing a safe, repeatable, and unbounded stream of interactions \cite{azoulay_adaptive_2021,deshmukh_developing_2025}. Second, it supports controlled comparison: rule-based, bandit-style, and RL policies can be evaluated under identical learner trajectories, isolating the contribution of the learning algorithm \cite{azoulay_adaptive_2021,axak_adaptive_nodate,weitekamp_tutorgym_2025}. Third, it enables systematic ablation, in particular contrasting emotion-aware and emotion-agnostic state representations under matched simulation conditions \cite{weitekamp_tutorgym_2025}.

The principal limitation of this evaluation strategy is the gap between simulated and real learners: a policy that performs well in simulation may behave differently when deployed on actual students, because real learner dynamics are richer and only partially observable \cite{urbaite_adaptive_2026}. This thesis therefore treats simulation results as evidence about algorithmic behaviour rather than as direct estimates of classroom performance.
